\documentclass[11pt]{article}

%% Standard packages
\usepackage[letterpaper,left=1in,right=1in,top=1in,bottom=1.1in,includefoot,footskip=28pt]{geometry}
\usepackage{amsmath, amssymb}
\usepackage{amsthm}
\usepackage{mathtools}
\usepackage{bm}
\usepackage{bbm}
\usepackage{graphicx}
\usepackage{booktabs}
\usepackage{hyperref}
\usepackage{cleveref}
\usepackage{algorithm}
\usepackage{algorithmic}
\usepackage{xcolor}
\usepackage{microtype}
\usepackage{natbib}
\usepackage{fancyhdr}
\usepackage{titling}
\providecommand{\theHalgorithm}{\arabic{algorithm}}
\providecommand{\theHfigure}{\arabic{figure}}

%% Notation macros
\newcommand{\R}{\mathbb{R}}
\newcommand{\N}{\mathbb{N}}
\newcommand{\bx}{\mathbf{x}}
\newcommand{\bX}{\mathbf{X}}
\newcommand{\bz}{\mathbf{z}}
\newcommand{\bZ}{\mathbf{Z}}
\newcommand{\bd}{\mathbf{d}}
\newcommand{\bD}{\mathbf{D}}
\newcommand{\bB}{\mathbf{B}}
\newcommand{\bI}{\mathbf{I}}
\newcommand{\bJ}{\mathbf{J}}
\newcommand{\bU}{\mathbf{U}}
\newcommand{\bV}{\mathbf{V}}
\newcommand{\bW}{\mathbf{W}}
\newcommand{\norm}[1]{\left\|#1\right\|}
\newcommand{\abs}[1]{\left|#1\right|}
\newcommand{\innerprod}[2]{\langle #1, #2 \rangle}
\newcommand{\dgeo}{d^{\mathrm{geo}}}
\newcommand{\demb}{d^{\mathrm{emb}}}
\newcommand{\dknn}{d^{k\text{-NN}}}
\newcommand{\Gk}{G_k}
\newcommand{\Stress}{\mathrm{Stress}}
\newcommand{\StressOne}{S_1}
\newcommand{\sgn}{\operatorname{sgn}}
\newcommand{\diag}{\operatorname{diag}}
\newcommand{\placeholderfigurebox}[2][0.9\linewidth]{%
  \fbox{\parbox[c][0.18\textheight][c]{#1}{\centering\footnotesize #2}}%
}
\DeclareMathOperator*{\argmin}{arg\,min}
\setlength{\droptitle}{-2.2em}
\setlength{\headheight}{14pt}
\pagestyle{plain}

\theoremstyle{definition}
\newtheorem{definition}{Definition}[section]
\theoremstyle{plain}
\newtheorem{proposition}{Proposition}[section]
\theoremstyle{remark}
\newtheorem{remark}{Remark}[section]

\newenvironment{keywords}{%
  \par\smallskip\noindent\textbf{Keywords: }\ignorespaces
}{\par\smallskip}

\newcommand{\paperaffiliationfooter}{%
  \parbox{0.92\textwidth}{\centering\footnotesize
  Center for Advanced Microbiome Research and Innovation (CAMRI), Institute for Genome Sciences, Department of Microbiology and Immunology, University of Maryland School of Medicine, Baltimore, Maryland, USA.}%
}

\fancypagestyle{firstpage}{%
  \fancyhf{}%
  \renewcommand{\headrulewidth}{0pt}%
  \renewcommand{\footrulewidth}{0pt}%
  \fancyfoot[C]{\paperaffiliationfooter}%
}

\providecommand{\manuscriptversion}{unversioned}
\providecommand{\manuscriptbuildnumber}{NA}
\providecommand{\manuscriptcommit}{unknown}
\providecommand{\manuscriptbuilddatetime}{unknown}
\InputIfFileExists{../build/manuscript_build_info.tex}{}{}

\begin{document}

\title{Geodesic Multidimensional Scaling:\\
       Self-Consistent Dimensionality Reduction\\ via Graph Geodesic Distances}

\author{Pawel Gajer \and Jacques Ravel}
\date{}

\maketitle
\thispagestyle{firstpage}
\begin{center}
\footnotesize
Version \texttt{\manuscriptversion} \quad | \quad Build
\texttt{\manuscriptbuildnumber} \quad | \quad Commit
\texttt{\manuscriptcommit}\\
Built \manuscriptbuilddatetime
\end{center}
\vspace{0.6em}

%% ============================================================
\begin{abstract}
Multidimensional scaling (MDS) seeks a low-dimensional embedding of
a set of points that faithfully reproduces a given pairwise dissimilarity
matrix.  Classical and metric MDS measure faithfulness via the discrepancy
between the input dissimilarities and the \emph{Euclidean} distances among
the embedded points.  Isomap improved upon MDS for curved-manifold data by
replacing Euclidean input dissimilarities with shortest-path distances on a
$k$-nearest-neighbor ($k$-NN) graph, thereby approximating intrinsic
manifold geodesics.  However, Isomap retains Euclidean distances on the
\emph{output} side, creating a fundamental asymmetry: input distances are
geodesic, but the fidelity criterion is Euclidean.  We argue that this
asymmetry is not merely an aesthetic inconsistency but a systematic source
of distortion whenever the data carry genuine graph or manifold structure.

We propose \emph{Geodesic MDS} (GMDS), a self-consistent formulation for
finite shortest-path metric spaces in which the output fidelity criterion
measures the embedded lengths of the same input-side geodesics that define
the input distances.  This separates two conceptually distinct steps that are
often conflated in practice: first, modeling a dataset by a weighted graph or
another shortest-path realization; second, embedding that geodesic metric
space in $\R^d$.  We define GMDS on a weighted graph together with a chosen
geodesic family, derive the corresponding geodesic stress and its gradient,
and discuss optimization strategies ranging from pure GMDS to regularized and
multiscale seeded variants.  The current implementation story is already
informative: on synthetic meshes the smooth target geometry may exist as an
exact GMDS solution while naive optimization still misses it, revealing that
the central practical problem is basin selection rather than formal
well-posedness alone.  We therefore present GMDS both as a mathematical
proposal and as an ongoing computational program whose current experiments
identify the most promising next steps toward a reliable working method.
\end{abstract}

\begin{keywords}
  multidimensional scaling, dimensionality reduction, geodesic distance,
  graph metric, $k$-nearest neighbor graph, Isomap, SMACOF, metric geometry,
  manifold learning, stress function
\end{keywords}

%% ============================================================
\section{Introduction}
\label{sec:introduction}

The central task of \emph{multidimensional scaling} (MDS) is to find a
configuration of $n$ points $\bZ = (\bz_1, \ldots, \bz_n) \in \R^{n \times
d}$ in a low-dimensional Euclidean space $\R^d$ such that the Euclidean
distances $\norm{\bz_i - \bz_j}$ are as close as possible to a given $n
\times n$ symmetric, non-negative dissimilarity matrix $\Delta = (\delta_{ij})$
\citep{torgerson1952multidimensional, kruskal1964multidimensional,
borg2005modern, cox2000multidimensional}.  The criterion of closeness is made
precise by a \emph{stress} function, and the embedding is obtained by
minimizing this stress.  MDS and its relatives are among the oldest and most
widely used tools in multivariate statistics, visualization, and machine
learning \citep{shepard1962analysis, sammon1969nonlinear,
vandermaaten2008tsne, mcinnes2018umap}.

When the input data are given as a matrix $\bX \in \R^{n \times m}$ of $n$
observations in a high-dimensional ambient space $\R^m$, the most common
choice is $\delta_{ij} = \norm{\bx_i - \bx_j}_{\R^m}$, the Euclidean
distance in the original space.  This is appropriate when the data are
believed to lie in a linear subspace of $\R^m$, so that Euclidean distances
faithfully reflect intrinsic structure.  However, when the data lie on or
near a nonlinear manifold embedded in $\R^m$, straight-line Euclidean
distances are poor proxies for the intrinsic distances that govern the
geometry of the data: two points near opposite ends of a curved surface
may be close in ambient Euclidean distance but far apart along the surface.

This limitation motivated \emph{Isomap}
\citep{tenenbaum2000global}, one of the most influential manifold learning
algorithms.  Isomap approximates the intrinsic geodesic distance between
points by constructing a $k$-nearest-neighbor ($k$-NN) graph on the data,
computing all-pairs shortest paths on this graph (weighted by Euclidean
edge lengths), and then feeding these graph-geodesic dissimilarities
$\delta_{ij} = \dknn(\bx_i, \bx_j)$ into classical MDS.  For data lying on
a smooth, sufficiently sampled Riemannian manifold, the graph-geodesic
distances converge to the true manifold geodesics as $n \to \infty$
\citep{bernstein2000graph, tenenbaum2000global}, and the resulting embedding
recovers the intrinsic geometry with high fidelity.

This paper argues that two independent issues should be kept distinct.  The
first is a \emph{modeling} issue: given a dataset $\bX$, how should one build
a weighted graph or other shortest-path realization that represents its
intrinsic geometry?  The second is an \emph{embedding} issue: given a
shortest-path metric space, how should one define an embedding criterion that
is self-consistent with that metric structure?  Geodesic MDS addresses the
second question directly and treats the first as an important but separate
construction layer.

\paragraph{The output-side inconsistency.}

Despite this advance, Isomap retains a fundamental asymmetry that has not
previously been identified as a source of systematic distortion.  On the
\emph{input} side, distances are geodesic---shortest paths through the data
graph.  On the \emph{output} side, the stress function measures faithfulness
via \emph{Euclidean} distances between embedded points.  This means that
Isomap optimizes an embedding in which the distance between $\bz_i$ and
$\bz_j$ is the straight-line distance $\norm{\bz_i - \bz_j}$, not the
embedded length of the corresponding input-side geodesic.

The comparison below summarizes the distinction between standard MDS,
Isomap, and the proposed Geodesic MDS in terms of the distance used on
the input and output sides of the embedding problem.

\begin{center}
  \small
  \begin{tabular}{@{}p{0.19\textwidth}p{0.34\textwidth}p{0.34\textwidth}@{}}
    \toprule
    Method & Input dissimilarity $\delta_{ij}$ & Output distance $d_{ij}$ \\
    \midrule
    Standard MDS & Euclidean $\norm{\cdot}_{\R^m}$ & Euclidean $\norm{\cdot}_{\R^d}$ \\
    Isomap & $d_{G_k}$ (graph geodesic) & Euclidean $\norm{\cdot}_{\R^d}$ \\
    Geodesic MDS & $d_{G_k}$ (graph geodesic) & embedded length of a fixed input geodesic \\
    \bottomrule
  \end{tabular}
\end{center}

For data living on a smooth manifold with slowly varying curvature this
asymmetry causes only a small error: nearby embedded points have
approximately equal Euclidean and geodesic distances.  But for data with
genuine graph or tree structure---where the topology enforces that all
paths between certain pairs of points must pass through bottleneck
vertices---the Euclidean output distance is a systematic \emph{underestimate}
of the embedded length of the relevant graph geodesic (by the triangle inequality, with
equality only when points are collinear in the embedding space).

Consider a concrete example: a star graph with center $c$ and three arm
tips $v_1, v_2, v_3$, each arm of unit length.  The graph distance between
$v_1$ and $v_2$ is $\delta(v_1, v_2) = 2$ (through $c$).  Any faithful
embedding must respect the topological constraint that the only path from
$v_1$ to $v_2$ passes through $c$.  Standard MDS and Isomap both try to
place $\bz_{v_1}$ and $\bz_{v_2}$ at Euclidean distance $2$ apart, which
is perfectly achievable in $\R^2$---but doing so forces the embedded arms
to be nearly collinear, misrepresenting the star topology.  The \emph{correct}
criterion is that the \emph{path-length distance along the embedded arms}
from $\bz_{v_1}$ to $\bz_{v_2}$ equals $\norm{\bz_{v_1} - \bz_c} +
\norm{\bz_c - \bz_{v_2}} = 2$, which is compatible with any angular
arrangement of the arms and correctly reflects the star structure.
Figure~\ref{fig:star_bias} is intended to make this contrast visually explicit.

\begin{figure}[t]
  \centering
  \includegraphics[width=0.98\linewidth]{../figures/fig_star_bias.pdf}
  \caption{Branching metrics expose the output-side inconsistency.  In a star graph, the graph distance between leaves is realized through the center.  Matching this distance with Euclidean output distances biases the embedding toward collinearity, whereas measuring output fidelity by the embedded lengths of the corresponding input geodesics preserves the correct topological structure.}
  \label{fig:star_bias}
\end{figure}

\paragraph{Contributions.}

This paper introduces \emph{Geodesic MDS}, a reformulation of the metric
MDS problem that is self-consistent with respect to shortest-path distance.
The main contributions of the present draft are as follows.

\begin{enumerate}
  \item \textbf{Conceptual separation.} We distinguish the step that maps a
  dataset to a weighted graph or shortest-path model from the GMDS problem
  proper, which is defined on that geodesic metric space once chosen.

  \item \textbf{General GMDS formulation.} We define GMDS on a weighted
  graph together with a chosen geodesic family $(G,\Gamma)$, thereby making
  the method explicit as a shortest-path-metric embedding method rather than
  only as a $k$-NN-data pipeline.

  \item \textbf{Optimization viewpoint.} We derive the geodesic stress,
  characterize its path-nonlocal gradient structure, and summarize the
  current optimization variants: pure GMDS, regularized GMDS, edge-based
  relaxations, and MISF-based multiscale schedules.

  \item \textbf{Current empirical understanding.} We summarize what the
  recent experiments now show: the smooth target often exists as a GMDS
  solution, but initialization and basin selection dominate practical
  success on the difficult synthetic families considered so far.

  \item \textbf{Research agenda.} We make explicit which parts of the story
  are already mathematically clear and which remain active computational
  problems, especially seeded multiscale optimization and truth-free seed
  selection.
\end{enumerate}

\paragraph{Relationship to existing work.}

Geodesic MDS is closely related to, but distinct from, several existing
methods.  \emph{Isomap} \citep{tenenbaum2000global} uses graph geodesics
on the input side only.  \emph{Kamada-Kawai graph drawing}
\citep{kamada1989algorithm} minimizes a stress between graph-theoretic
shortest-path distances and Euclidean drawing distances, but is restricted
to 2D visualization and is not formulated as a dimensionality reduction
method with explicit input/output distance comparison.  \emph{Metric graph
embedding} theory \citep{bourgain1985lipschitz,
linial1994geometry} studies distortion bounds for embedding graph metrics
into $\ell_p$ spaces but does not define a stress function or optimization
criterion.  \emph{Diffusion maps} \citep{coifman2006diffusion} define a
Markov-chain-based distance that is intrinsic to the data graph but again
uses Euclidean output distances.  To our knowledge, the self-consistent
formulation proposed here---graph geodesics on the input side and
embedded lengths of the corresponding fixed input geodesics on the output
side---has not been previously proposed.  The present paper frames that
proposal first at the level of shortest-path metric spaces and only then
returns to data-driven graph constructions such as $k$-NN models.

\paragraph{Organization.}

The remainder of the paper is organized as follows.
\cref{sec:background_v4} revisits MDS, Isomap, and the shortest-path
viewpoint needed to motivate the self-consistency principle.
\cref{sec:geodesic_mds_v4} then defines GMDS on shortest-path metric spaces.
\cref{sec:optimization_v4} turns to optimization and practical variants,
including the current multiscale implementations.  \cref{sec:experiments_v4}
summarizes what the current implementation and experiments actually teach us,
and \cref{sec:discussion_v4} closes with the next steps needed to move the
method toward a more reliable working implementation.  Technical background on
classical MDS, solver derivations, and metric-geometric context is collected in
the appendices so that the main text can follow the scientific story more
directly.

%% ============================================================
\section{From MDS to Geodesic Self-Consistency}
\label{sec:background_v4}
\label{sec:background}

This section isolates the conceptual transition from Euclidean stress-based
embedding to geodesically self-consistent embedding.  The key organizational
point is that two questions should be kept distinct.  The first is a modeling
question: given a dataset, how should we build a weighted graph or another
shortest-path realization?  The second is an embedding question: once such a
shortest-path metric space has been chosen, how should its fidelity be measured
in a low-dimensional embedding?  GMDS is about the second question directly.

\subsection{Classical and Metric MDS}
\label{sec:classical_mds}
\label{sec:metric_mds}

Classical and metric MDS take as input a symmetric dissimilarity matrix
$\Delta=(\delta_{ij})$ and seek an embedding $\bZ\in\R^{n\times d}$ for which
the Euclidean output distances $\norm{\bz_i-\bz_j}$ reproduce those
dissimilarities as closely as possible
\citep{torgerson1952multidimensional, young1938discussion,
kruskal1964multidimensional, borg2005modern}.  When $\Delta$ is Euclidean,
classical MDS solves the problem in closed form through the double-centered
Gram matrix
\begin{equation}
  \bB = -\tfrac{1}{2} \bJ \bD^{(2)} \bJ,
  \label{eq:gram}
\end{equation}
with $\bD^{(2)}_{ij}=\delta_{ij}^2$ and
$\bJ=\bI_n-n^{-1}\mathbf{1}\mathbf{1}^\top$.  When $\Delta$ is not Euclidean,
metric MDS instead minimizes a stress criterion such as
\begin{equation}
  \StressOne(\bZ) = \sqrt{
    \frac{\displaystyle\sum_{i < j}
          \bigl(d_{ij}(\bZ) - \delta_{ij}\bigr)^2}
         {\displaystyle\sum_{i < j} d_{ij}(\bZ)^2}
  },
  \label{eq:stress1}
\end{equation}
or its weighted analogue
\begin{equation}
  \sigma_w(\bZ)
  = \sum_{i < j} w_{ij}\,\bigl(d_{ij}(\bZ) - \delta_{ij}\bigr)^2.
  \label{eq:weighted_stress}
\end{equation}
For the purposes of this paper, the important fact is simple: the output-side
quantity being matched is always the Euclidean distance in the embedding space.
Additional background on classical, metric, non-metric, and SMACOF-style MDS
is deferred to \cref{app:mds_background}.

\subsection{Isomap and the Output-Side Inconsistency}
\label{sec:isomap}

Isomap \citep{tenenbaum2000global} changes the input side of the problem by
replacing ambient Euclidean dissimilarities with shortest-path distances on a
weighted neighborhood graph.  Given $\bX\in\R^{n\times m}$, Isomap:
\begin{enumerate}
  \item constructs a weighted graph, typically a $k$-NN graph;
  \item computes all-pairs shortest-path distances on that graph; and
  \item applies classical MDS to that new distance matrix.
\end{enumerate}
This is a major conceptual advance because it recognizes that ambient Euclidean
distance is often the wrong \emph{input} geometry.  But Isomap still evaluates
the embedding by Euclidean output distances $\norm{\bz_i-\bz_j}$.  For smooth
manifolds this may be a small approximation; for graph-like, branching, or
bottlenecked metrics it creates a systematic mismatch between the type of
distance used on the input side and the type of distance used on the output
side.

\subsection{From Data to a Weighted Graph or Shortest-Path Model}
\label{sec:knn_graph}

Once the modeling and embedding questions are separated, the role of the graph
becomes clearer.  A dataset may be converted into a weighted graph or another
finite shortest-path realization in many ways: $k$-NN graphs, radius graphs,
minimum realizations, application-specific weighted networks, or a graph that
is already given a priori.  GMDS does not prescribe which construction is
correct.  Instead it starts from a chosen shortest-path metric space and asks
for a self-consistent embedding of that space.

For data analysis, the neighborhood parameter $k$ remains important because it
controls the geometric scale at which the graph path metric departs from the
ambient Euclidean one.  Small $k$ emphasizes sparse, graph-like structure;
large $k$ drives the model back toward the ambient Euclidean regime.  This is
illustrated in \cref{fig:k_regimes}.

\begin{figure}[t]
  \centering
  \includegraphics[width=0.98\linewidth]{../figures/fig_k_regimes.pdf}
  \caption{The neighborhood parameter $k$ acts as a geometric scale parameter.  The same horseshoe-shaped point set is shown for three neighborhood sizes, with the same endpoint pair $p,q$ highlighted in each panel.  For small $k$, the graph is sparse and nearly tree-like; for intermediate $k$, shortest paths still follow the manifold and preserve local geodesic structure; and for large $k$, dense shortcut edges appear and the metric approaches the ambient Euclidean regime.}
  \label{fig:k_regimes}
\end{figure}

\subsection{The Main Transition}

The transition to GMDS can therefore be stated compactly.  Standard MDS uses
Euclidean dissimilarities on input and Euclidean distances on output.  Isomap
replaces only the input distance by a graph-geodesic distance.  GMDS completes
the transition by comparing geodesic quantities on both sides: a shortest-path
metric on the input side and the embedded lengths of the same chosen geodesics
on the output side.  The next section makes that definition explicit on finite
shortest-path metric spaces.

%% ============================================================
\section{Geodesic MDS on Shortest-Path Metric Spaces}
\label{sec:geodesic_mds_v4}
\label{sec:geodesic_mds}

We now leave the dataset-construction step aside and state the GMDS problem in
its own natural category: finite shortest-path metric spaces, represented in
practice by weighted graphs together with a chosen family of geodesics.

\subsection{Self-Consistency Principle}
\label{sec:self_consistency}

We begin by articulating the principle that motivates the proposed method.

\begin{definition}[Self-consistent embedding]
  \label{def:self_consistency}
  Let $(V, d_G)$ be a finite shortest-path metric space realized by a
  connected weighted graph $G = (V, E, w)$.  For each pair $(i,j)$, choose a
  canonical $d_G$-geodesic $\gamma_{ij}$ in $G$.  An embedding
  $f: V \to \R^d$ is
  \emph{self-consistent} with respect to $(G, \Gamma)$, where
  $\Gamma = \{\gamma_{ij}\}_{i<j}$, if the fidelity of the embedding is
  measured by the path length of the image $f(\gamma_{ij})$---i.e., by the
  sum of Euclidean edge lengths along that same input-side geodesic.
\end{definition}

Standard metric MDS and Isomap are not self-consistent in this sense:
they measure fidelity by Euclidean distances $\norm{f(i) - f(j)}$, which
are at most equal to any path-length distance through intermediate vertices.

The self-consistency principle has a natural interpretation:
\emph{the distance used to evaluate the embedding should be the same
type of distance, along the same family of geodesics, used to construct
the input dissimilarities.}
If input dissimilarities are graph geodesics, output distances should be
measured along those same graph geodesics after embedding.  This is not
merely a philosophical preference: the
analysis in \cref{sec:systematic_bias} below shows that using Euclidean
output distances introduces a systematic downward bias in the estimated
distances for topologically constrained pairs.

\subsection{Systematic Bias of Euclidean Output Distances}
\label{sec:systematic_bias}

\begin{proposition}[Triangle inequality bias]
  \label{prop:bias}
  Let $\bZ = (\bz_1, \ldots, \bz_n) \in \R^{n \times d}$ be any embedding,
  and let $G = (V, E, w)$ be a weighted graph with vertex set $V = \{1,
  \ldots, n\}$.  For any pair $(i,j)$ with chosen shortest path
  $i = v_0, v_1, \ldots, v_\ell = j$ in $G$, define the embedded length of
  that geodesic by
  \begin{equation}
    \demb_G(i,j;\bZ)
    = \sum_{s=0}^{\ell-1} \norm{\bz_{v_s} - \bz_{v_{s+1}}}.
    \label{eq:emb_geo}
  \end{equation}
  Then
  \begin{equation}
    \norm{\bz_i - \bz_j} \;\leq\; \demb_G(i,j;\bZ),
    \label{eq:triangle_bias}
  \end{equation}
  with equality if and only if all embedded path vertices
  $\bz_{v_0}, \ldots, \bz_{v_\ell}$ are collinear in $\R^d$ in the order
  of the path.
\end{proposition}

\begin{proof}
  By the triangle inequality applied $\ell$ times:
  $\norm{\bz_i - \bz_j} = \norm{\sum_{s} (\bz_{v_{s+1}} - \bz_{v_s})}
  \leq \sum_s \norm{\bz_{v_{s+1}} - \bz_{v_s}} = \demb_G(i,j;\bZ)$.
  Equality holds iff all difference vectors $\bz_{v_{s+1}} - \bz_{v_s}$
  are non-negative scalar multiples of a single unit vector, i.e., the
  points are collinear in the stated order.
\end{proof}

\Cref{prop:bias} shows that standard MDS and Isomap systematically
underestimate embedded geodesic lengths for all topologically
non-degenerate pairs (those whose shortest path has $\ell \geq 2$ and
does not embed as a straight line).  When the input dissimilarity
$\delta_{ij}$ is a graph geodesic, minimizing Euclidean-output stress
creates pressure to satisfy $\norm{\bz_i - \bz_j} \approx \delta_{ij}$,
which in turn forces embedded paths to be approximately collinear.
For tree-structured data this collapses the tree into a lower-dimensional
structure than warranted; for star graphs it forces nearly collinear arm
arrangement regardless of the dimension $d$.

\subsection{Geodesic Stress on $(G,\Gamma)$}
\label{sec:geodesic_stress}

We now define the proposed method.

\paragraph{Setup.}
Let $G = (V,E,w)$ be a connected weighted graph with vertex set
$V = \{1,\ldots,n\}$ and positive edge weights $w_{uv} > 0$.  Write
$d_G$ for the weighted shortest-path metric on $V$.  For each pair
$i < j$, choose a canonical $d_G$-geodesic
\begin{equation}
  \gamma_{ij} = (v_0 = i, v_1, \ldots, v_L = j),
  \label{eq:canonical_geodesic}
\end{equation}
breaking ties deterministically when multiple shortest paths exist.  The
collection $\Gamma = \{\gamma_{ij} : i < j\}$ is part of the data of the
method.  In this sense, Geodesic MDS is defined on the pair $(G,\Gamma)$,
not on $G$ alone.

An embedding is a configuration matrix $\bZ = (\bz_1, \ldots, \bz_n)^\top
\in \R^{n \times d}$.  For fixed $(G,\Gamma)$, the \emph{embedded graph}
$G(\bZ)$ is the same graph $G$ with edge lengths replaced by embedded
Euclidean distances.  The \emph{embedded length of the chosen input
geodesic} is
\begin{equation}
  \demb_{\Gamma}(i,j;\bZ)
  = \sum_{s=0}^{L-1} \norm{\bz_{v_s} - \bz_{v_{s+1}}},
  \label{eq:emb_geodesic}
\end{equation}
where $\gamma_{ij} = (v_0, \ldots, v_L)$.

\begin{definition}[Raw geodesic stress]
  \label{def:raw_geodesic_stress}
  The \emph{raw geodesic stress} of an embedding $\bZ \in \R^{n \times d}$
  with respect to $(G,\Gamma)$ is
  \begin{equation}
    \Stress^{\mathrm{geo}}_{\Gamma}(\bZ; G)
    = \sum_{i < j}
      \bigl(\demb_{\Gamma}(i,j;\bZ) - d_G(i,j)\bigr)^2.
    \label{eq:raw_geo_stress}
  \end{equation}
\end{definition}

\begin{definition}[Geodesic Stress-1]
  \label{def:geodesic_stress}
  The normalized \emph{geodesic Stress-1} of $\bZ$ with respect to
  $(G,\Gamma)$ is
  \begin{equation}
    {\StressOne}_{\Gamma}^{\mathrm{geo}}(\bZ; G)
    = \sqrt{
        \frac{\Stress^{\mathrm{geo}}_{\Gamma}(\bZ; G)}
             {\displaystyle\sum_{i < j} d_G(i,j)^2}
      }.
    \label{eq:geo_stress}
  \end{equation}
\end{definition}

\begin{definition}[Geodesic MDS embedding]
  \label{def:gmds_embedding}
  A \emph{Geodesic MDS embedding} of $(G,\Gamma)$ in $\R^d$ is any
  minimizer
  \begin{equation}
    \hat{\bZ}_{\Gamma,d}(G)
    \in \argmin_{\bZ \in \R^{n \times d}}
    \Stress^{\mathrm{geo}}_{\Gamma}(\bZ; G).
    \label{eq:geo_mds}
  \end{equation}
  Equivalently, since the denominator in \cref{eq:geo_stress} is constant
  in $\bZ$, one may minimize ${\StressOne}_{\Gamma}^{\mathrm{geo}}(\bZ; G)$.
\end{definition}

\paragraph{Specialization to $k$-NN graphs.}
When the input is a data matrix $\bX \in \R^{n \times m}$, we recover the
construction used throughout this paper by fixing a neighborhood parameter
$k \in \{1,\ldots,n-1\}$ and building the symmetric $k$-NN graph
$\Gk = (V, E_k)$ with
\begin{equation}
  E_k = \{(i,j) : j \in \mathcal{N}_k(i) \text{ or } i \in \mathcal{N}_k(j)\},
  \label{eq:knn_edges}
\end{equation}
where $\mathcal{N}_k(i)$ denotes the $k$ nearest neighbors of $i$ in
$\R^m$.  Assign weights $w_{ij} = \norm{\bx_i - \bx_j}_{\R^m}$ for
$(i,j) \in E_k$, compute the shortest-path metric
$d_{\Gk}(i,j) = \delta_{ij}^{(k)}$, and choose a canonical shortest-path
family $\Gamma_k = \{\gamma_{ij}^{(k)} : i < j\}$.  The resulting Geodesic
MDS objective is
\begin{equation}
  \Stress^{\mathrm{geo}}_{\Gamma_k}(\bZ; \Gk)
  = \sum_{i < j}
    \bigl(\demb_{\Gamma_k}(i,j;\bZ) - \delta^{(k)}_{ij}\bigr)^2.
  \label{eq:knn_raw_geo_stress}
\end{equation}

\begin{remark}[Normalization]
  We normalize by the constant $\sum_{i<j} d_G(i,j)^2$ rather than by an
  embedded denominator such as $\sum_{i<j} \demb_{\Gamma}(i,j;\bZ)^2$.
  This keeps ${\StressOne}_{\Gamma}^{\mathrm{geo}}$ scale-free while
  leaving the optimizer identical to the raw numerator
  \cref{eq:raw_geo_stress}.  A Kruskal-style denominator depending on
  $\bZ$ is still useful for reporting, but it changes the optimization
  target.
\end{remark}

\begin{remark}[Fixed topology and fixed geodesics]
  The graph topology $E$ and the chosen geodesic family $\Gamma$ are fixed
  by the input metric data and are not part of the optimization.  Only the
  vertex positions $\bz_1,\ldots,\bz_n$ are free.  This is a key design
  choice: it ensures that the relevant routes are determined by the input
  geometry and that the optimization preserves that chosen geodesic
  structure after embedding.
\end{remark}
Figure~\ref{fig:geodesic_stress_definition} should make the fixed-geodesic,
two-metric comparison in \cref{eq:emb_geodesic,eq:geo_stress} easy to see.

\begin{figure}[t]
  \centering
  \includegraphics[width=0.98\linewidth]{../figures/fig_geodesic_stress_definition.pdf}
  \caption{Definition of geodesic stress.  Left: the input
  dissimilarity $d_G(i,j)$ is the shortest-path distance between $i$ and
  $j$ on the fixed weighted graph $G$; in the data-analytic specialization,
  $G = \Gk$ is a $k$-NN graph.  Right: the output quantity
  $\demb_{\Gamma}(i,j;\bZ)$ is the length of the canonical shortest path
  chosen in the input graph, measured after the vertex positions move in
  $\R^d$.  The route is fixed by the input geometry; only the embedded edge
  lengths change during optimization.}
  \label{fig:geodesic_stress_definition}
\end{figure}


\begin{remark}[Nonuniqueness is intrinsic, not a defect]
  The objective \cref{eq:raw_geo_stress} is a fidelity criterion, not an
  aesthetics criterion, so its minimizers need not be unique.  This is a
  feature of the definition: multiple embeddings can realize the same graph
  geodesic data equally well.
\end{remark}

For example, for a path graph $v_0 - v_1 - v_2$ with any fixed edge lengths,
every embedding in $\R^2$ that preserves those two edge lengths is a global
minimum of \cref{eq:raw_geo_stress}; one obtains an infinite family of minima
by fixing the segment $[v_0,v_1]$ and rotating $[v_1,v_2]$ around $v_1$.
Likewise, if $v_0,v_1,v_2$ form one triangle and $v_1,v_2,v_3$ form another
triangle sharing the edge $[v_1,v_2]$, then both the non-overlapping and the
overlapping realizations are global minima whenever all edge lengths are
preserved.  \Cref{fig:nonuniqueness_examples} illustrates these two basic
sources of nonuniqueness.

\begin{figure}[t]
  \centering
  \includegraphics[width=0.98\linewidth]{../figures/fig_nonuniqueness_examples.pdf}
  \caption{Two sources of nonuniqueness for geodesic-stress minimizers.
  Left: a three-vertex path has a one-parameter family of exact realizations
  obtained by rotating the second edge while preserving both edge lengths.
  Right: two triangles sharing an edge admit distinct exact realizations,
  including a non-overlapping embedding and an overlapping embedding.  In both
  examples the raw geodesic stress is zero because all relevant graph-geodesic
  lengths are preserved.}
  \label{fig:nonuniqueness_examples}
\end{figure}

\paragraph{A diagnostic lesson.}
The diagnostic experiments reported separately in this project show that the
all-pairs objective \cref{eq:raw_geo_stress} can admit visually pathological
minimizers on simple synthetic graphs.  On orthogonal grids, graph geodesics
systematically exceed straight-line Euclidean distances, so the gradient of the
all-pairs geodesic stress can push the embedding toward stretched or folded
configurations.  This does not invalidate the definition of GMDS: it shows that
regularization is a separate modeling choice that may be needed when one wants
both fidelity to graph geodesics and a visually well-spread realization.

\paragraph{Regularized GMDS.}
To incorporate such modeling choices explicitly, introduce a Euclidean
dispersion regularizer
\begin{equation}
  R(\bZ) = \sum_{i < j} f_r\bigl(\norm{\bz_i - \bz_j}\bigr),
  \label{eq:repulsive_regularizer}
\end{equation}
where $f_r : (0,\infty) \to (0,\infty)$ is a decreasing function.  A
concrete choice used widely in force-directed graph drawing is the Coulomb-type
penalty
\begin{equation}
  f_r(r) = \frac{1}{r^2 + \varepsilon_r^2},
  \label{eq:coulomb_regularizer}
\end{equation}
with a small $\varepsilon_r > 0$ to prevent singularities.  The
self-consistency principle applies to the fidelity term
\cref{eq:raw_geo_stress}; the role of $R$ is different.  It is a Euclidean
regularizer that encourages spread and discourages collapse, not a substitute
for the graph-geodesic mismatch term itself.  A standard scalarized objective
is therefore
\begin{equation}
  F_\lambda(\bZ; G,\Gamma)
  = \Stress^{\mathrm{geo}}_{\Gamma}(\bZ; G) + \lambda R(\bZ),
  \qquad \lambda > 0.
  \label{eq:regularized_gmds}
\end{equation}
A \emph{regularized GMDS embedding} is a local minimizer of
\cref{eq:regularized_gmds} for a specified $\lambda$.  The value of
$\lambda$ is not universal; in practice it should be reported explicitly and
chosen by a fidelity-dispersion trade-off study, continuation scheme, or
task-specific validation criterion.

\paragraph{An edge-based relaxation.}
If exact minimization of \cref{eq:raw_geo_stress} is computationally
prohibitive, one can replace it by the edge-only functional
\begin{equation}
  E_K(\bZ)
  = \sum_{(i,j) \in E} k_{ij}\bigl(\norm{\bz_i - \bz_j} - w_{ij}\bigr)^2,
  \label{eq:edge_relaxation}
\end{equation}
because for an edge $(i,j)$ the chosen graph geodesic is the edge itself and
thus $\demb_{\Gamma}(i,j;\bZ) = \norm{\bz_i - \bz_j}$.  This is no longer the
geodesic stress: it is a graph-drawing-style relaxation that penalizes only
single-hop distortion, and is therefore closer to an edge-restricted
Kamada-Kawai objective \citep{kamada1989algorithm, fruchterman1991graph}
than to the full all-pairs GMDS criterion.

The relation between \cref{eq:raw_geo_stress,eq:edge_relaxation} is easiest to
see in edge-length coordinates.  Let $\ell(\bZ) \in \R^{|E|}$ be the vector of
embedded edge lengths, let $w \in \R^{|E|}$ collect the graph edge weights, and
let $A_{\Gamma}$ be the path-edge incidence matrix whose $(i,j)$-row indicates
which edges lie on the chosen geodesic $\gamma_{ij}$.  Then
\begin{equation}
  \Stress^{\mathrm{geo}}_{\Gamma}(\bZ; G)
  = \norm{A_{\Gamma}\bigl(\ell(\bZ) - w\bigr)}_2^2
  = \bigl(\ell(\bZ)-w\bigr)^\top
    A_{\Gamma}^\top A_{\Gamma}
    \bigl(\ell(\bZ)-w\bigr),
  \label{eq:path_incidence_stress}
\end{equation}
whereas
\begin{equation}
  E_K(\bZ)
  = \bigl(\ell(\bZ)-w\bigr)^\top K \bigl(\ell(\bZ)-w\bigr),
  \label{eq:edge_relaxation_matrix}
\end{equation}
with $K = \diag(k_{ij})$ diagonal.  Thus $E_K$ drops the off-diagonal
path-coupling terms in $A_{\Gamma}^\top A_{\Gamma}$ and therefore does not
control multi-hop geodesic distortion directly.  What is gained is a much
cheaper objective; what is lost is explicit control of long-range path-length
errors.

If $E_K(\bZ) = 0$, then every edge length is realized exactly, hence every
chosen graph-geodesic length is realized exactly as well, and
\cref{eq:raw_geo_stress} also vanishes.  In the approximate regime, however,
the two objectives are different and generally have different local minima.
For this reason we treat \cref{eq:edge_relaxation} as a relaxation of GMDS,
not as the definition of GMDS itself.  Whether this relaxation also avoids
some of the pathologies of the full all-pairs objective is an empirical
question, not a formal consequence of the definition.

The weights $k_{ij}$ in \cref{eq:edge_relaxation} are user-chosen tuning
parameters.  A simple default is $k_{ij} \equiv 1$.  If one wants to emphasize
edge lengths that occur frequently in the graph, one may set
$k_{ij} = \hat f_{\mathrm{edge}}(w_{ij})$, where
$\hat f_{\mathrm{edge}}$ is a kernel density estimate of the empirical
distribution of edge lengths.  Other choices, such as weights proportional to
geodesic-usage counts, are also possible and should be justified by the
application.

\subsection{Basic Properties of GMDS}
\label{sec:properties}

\paragraph{Self-consistency.}
Geodesic MDS satisfies \cref{def:self_consistency} by construction:
both input dissimilarities and output distances are measured on the same
fixed family of geodesics in $(G,\Gamma)$ (the former by the graph
metric $d_G$, the latter by the embedded path lengths
$\demb_{\Gamma}(i,j;\bZ)$).

\paragraph{Recovery of standard MDS at $k = n-1$.}
When the input graph is the $k$-NN specialization and $k = n-1$, every
pair is connected by a direct edge, so
$\demb_{\Gamma_{n-1}}(i,j;\bZ) = \norm{\bz_i - \bz_j}$ and
$\delta^{(n-1)}_{ij} = \norm{\bx_i - \bx_j}$.  Geodesic stress reduces to
standard metric Stress-1, and Geodesic MDS recovers standard metric MDS.

\paragraph{Tree metric at small $k$.}
For $k = 1$, $\Gk$ is (generically) a sparse forest; with a suitable
connectivity correction (e.g., taking the minimum spanning tree as
$G_1$) the graph is a tree.  The tree-path metric is then a tree metric in
the sense of \citet{dress1984trees}, and Geodesic MDS seeks an
embedding that faithfully represents the tree structure.  Tree metrics
embed isometrically into $\R^1$ (the real line) with zero stress;
thus for tree-structured data Geodesic MDS with $k=1$ achieves
perfect stress in $d=1$.

\paragraph{Interpolation and the role of $k$.}
As $k$ increases from 1 to $n-1$, more edges are added to $\Gk$,
$\delta^{(k)}_{ij}$ decreases toward $\norm{\bx_i - \bx_j}$, and
$\demb_{\Gamma_k}(i,j;\bZ)$ decreases toward $\norm{\bz_i - \bz_j}$.  The
geodesic stress thus interpolates between a tree-metric embedding at small
$k$ and standard MDS at large $k$.  The parameter $k$ controls the
geometric scale at which geodesic structure is enforced, analogous to the
bandwidth parameter in kernel methods.

\paragraph{Geodesic stress is never smaller than Euclidean stress.}
By \cref{prop:bias}, $\demb_{\Gamma_k}(i,j;\bZ) \geq \norm{\bz_i - \bz_j}$
for all $i,j$ and all $\bZ$.  If in addition $\delta^{(k)}_{ij} \geq
\norm{\bx_i - \bx_j}$ (which holds by the triangle inequality on graph
paths), then for the same configuration $\bZ$ the geodesic stress is
larger than the Euclidean stress.  The optimal geodesic configuration
$\hat{\bZ}_{\Gamma_k,d}(\Gk)$ is not in general the same as the optimal
Euclidean configuration.

\section{Optimization and Practical Variants}
\label{sec:optimization_v4}

The preceding section defines GMDS as a shortest-path-metric embedding
criterion.  We now turn to the current optimization story: what can already be
implemented, what serves only as a surrogate, and which multiscale seeded
variants currently offer the most promising practical behavior.

\subsection{Why Optimization is Difficult}
\label{sec:gradient}

The practical difficulty of GMDS is not only that the objective is non-convex.
It is also that the stress is \emph{path-nonlocal}.  Once a geodesic family
$\Gamma$ has been fixed, the contribution of a vertex depends on every endpoint
pair whose chosen input geodesic passes through that vertex.  The optimization
therefore couples local coordinate changes to global route structure, and the
resulting landscape is highly sensitive to initialization.  This is exactly why
the current implementation program has had to study basin selection, seeded
multiscale schedules, and lower-level placement mechanisms rather than treating
GMDS as ``just another stress solver.''  The full gradient formulas and a
majorization-based optimization viewpoint are recorded in
\cref{app:gradient_majorization}.

\begin{figure}[t]
  \centering
  \begin{minipage}[t]{0.82\textwidth}
    \centering
    \includegraphics[width=\linewidth]{../figures/fig_gradient_nonlocality_paths.pdf}
  \end{minipage}

  \caption{Geodesic stress has non-local gradients.  A vertex receives contributions only from endpoint pairs whose chosen input geodesic passes through it, so the gradient depends on global route structure rather than only on immediate graph neighbors.}
  \label{fig:gradient_nonlocality}
\end{figure}

\subsection{Pure GMDS, Regularized GMDS, and Surrogates}
\label{sec:geo_smacof}

In the current implementation program it is useful to distinguish three layers.
First, \emph{pure GMDS} minimizes the fixed-geodesic stress itself.  Second,
\emph{regularized GMDS} adds auxiliary terms such as anchor penalties,
graph-distance-aware repulsion, or edge-spring terms in order to steer the
optimization toward better basins.  Third, \emph{surrogates} such as
edge-relaxation objectives are used not as definitions of GMDS, but as basin
finders or controls in experiments.

The present paper does not claim that one particular solver has already solved
the optimization problem.  Instead, the computational contribution is a working
family of solvers whose behavior can now be studied scientifically.  The pure
objective is implemented in compiled C++ in the \texttt{grip} package; the
regularized variants share the same path-cache machinery; and the current
multiscale pipeline reuses MISF, GRIP, Weighted GRIP, and weighted KK as
basin-finding or placement mechanisms.  The detailed gradient formulas and a
path-majorization viewpoint are given in \cref{app:gradient_majorization}.

\begin{figure}[t]
  \centering
  \includegraphics[width=0.98\linewidth]{../figures/fig_geodesic_mds_workflow.pdf}
  \caption{High-level GMDS workflow.  A canonical family of shortest paths is chosen once from the input graph.  During optimization, the embedded lengths of those fixed geodesics are reevaluated repeatedly, and the resulting residuals drive whichever solver schedule is being used: pure GMDS, regularized GMDS, or a multiscale seeded schedule.}
  \label{fig:geodesic_mds_workflow}
\end{figure}

\subsection{Algorithmic Summary}

We summarize Geodesic MDS as \cref{alg:geodesic_mds}.

\begin{algorithm}
\caption{Geodesic MDS}
\label{alg:geodesic_mds}
\begin{algorithmic}[1]
  \REQUIRE Connected weighted graph $G=(V,E,w)$, chosen geodesic family
           $\Gamma$, target dimension $d$, convergence tolerance
           $\varepsilon>0$, maximum iterations $T$.
  \ENSURE Embedding $\bZ \in \R^{n \times d}$.
  \\
  \textit{// Step 1: Prepare the shortest-path model}
  \STATE Compute and cache the target graph-geodesic distances $d_G(i,j)$.
  \STATE Choose and store the fixed geodesic family
         $\Gamma = \{\gamma_{ij}: i<j\}$.
  \\
  \textit{// Step 2: Initialize}
  \STATE Initialize $\bZ^{(0)}$ from a seed pool
         (e.g.\ classical MDS, KK, GRIP, Weighted GRIP, or a multiscale scaffold).
  \STATE $t \leftarrow 0$.
  \\
  \textit{// Step 3: Optimize}
  \REPEAT
    \STATE Evaluate the embedded path lengths
           $\demb_{\Gamma}(i,j;\bZ^{(t)})$ along the stored geodesics.
    \STATE Update $\bZ^{(t)}$ using the current solver schedule
           (pure GMDS, regularized GMDS, or multiscale seeded refinement).
    \STATE $t \leftarrow t + 1$.
  \UNTIL{$\abs{\sigma^{\mathrm{geo}}(\bZ^{(t)}) - \sigma^{\mathrm{geo}}(\bZ^{(t-1)})} < \varepsilon$
         or $t \geq T$.}
  \RETURN $\bZ^{(t)}$.
\end{algorithmic}
\end{algorithm}

\paragraph{Initialization.}
Classical MDS applied to graph-geodesic distances remains an important baseline
seed, but the recent experiments show that it is not a universal best choice.
On highly regular meshes, alternatives such as KK, GRIP, Weighted GRIP, or
multiscale seeded scaffolds can lead to substantially better final GMDS
embeddings.  In practice, classical MDS is therefore best treated as one member
of an initializer pool rather than as a single default.

\paragraph{Multiple restarts.}
Since the geodesic stress landscape is non-convex, multiple restarts are
advisable.  The current implementation suggests that the most useful restarts
are often \emph{structured} rather than purely random: spectral seeds, KK
seeds, GRIP-family seeds, and multiscale coarse scaffolds explore genuinely
different basins, whereas small random perturbations of a single seed often do
not.

\paragraph{Selection of $k$.}
When the shortest-path model is itself constructed from data, the parameter $k$
controls the tradeoff between graph-like structure (small $k$) and ambient
Euclidean structure (large $k$).  We therefore view $k$ selection as part of
the modeling step rather than part of the GMDS definition.  In practice it can
be chosen by held-out geodesic stress, by stability under resampling, or by
geometric prior knowledge about the data.

\subsection{A MISF-based Multiscale Modification of GMDS}
\label{sec:misf_gmds}

The experiments of \cref{sec:properties} show that a global classical-MDS
initialization can place GMDS into an unfavorable basin, especially on highly
regular meshes.  A natural alternative is to borrow the
\emph{maximal-independent-set filtration} (MISF) strategy introduced for large
force-directed layouts by \citet{gajer2004multidimensional,gajer2002grip} and
replace the force-directed objective at each scale by the geodesic objective
developed above.  The resulting method is still a fixed-geodesic GMDS method:
the fidelity term remains \cref{eq:raw_geo_stress}, but the initialization and
the subsequent optimization are carried out in a coarse-to-fine manner.

\paragraph{Filtration.}
Given a connected weighted graph $G=(V,E,w)$, let
\[
  V = V_0 \supset V_1 \supset \cdots \supset V_L
\]
be a maximal-independent-set filtration in the sense of
\citet{gajer2004multidimensional,gajer2002grip}: each $V_{\ell+1}$ is a maximal
subset of $V_\ell$ whose vertices are separated in the graph by a prescribed
minimum graph distance, so that the levels become progressively coarser.
Intuitively, $V_L$ is a sparse sample of the original graph, while $V_0=V$ is
the full vertex set.  The filtration is shown on the underlying $2$D mesh
geometry in \cref{fig:misf_filtration_cases}.  For the regular $12\times 12$
paraboloid the extracted level sizes are $24 \rightarrow 50 \rightarrow 144$;
for the regular $15\times 15$ paraboloid they are
$24 \rightarrow 28 \rightarrow 84 \rightarrow 225$; and for the irregular
rectangle paraboloid they are $24 \rightarrow 32 \rightarrow 81 \rightarrow
225$.

\paragraph{Top level.}
Instead of initializing the whole graph by classical MDS, we first solve a
\emph{pure} GMDS problem only on the coarsest set $V_L$.  Concretely, we form
the complete graph on $V_L$ with edge weights equal to the graph-geodesic
distances $d_G(i,j)$ measured in the original graph $G$, choose a fixed
geodesic family $\Gamma_L$ on this induced metric graph, and minimize the raw
geodesic stress on $V_L$.  This coarse solve is the \emph{initial embedding} of
the MISF pipeline.  The middle panel in the top row of
\cref{fig:misf_stages_regular_12x12,fig:misf_stages_regular_15x15,fig:misf_stages_irregular_15x15}
shows this initial coarse embedding.  It is not a random seed: it is already
the result of a restartable pure-GMDS solve on $V_L$.  The right panel in the
same row shows that this coarse embedding is later revisited by the sparse
refinement stage.

\paragraph{Geodesic insertion of finer levels.}
Once $V_{\ell+1}$ has been embedded, each new vertex $v \in V_\ell \setminus
V_{\ell+1}$ is placed with respect to a small anchor set
$A_\ell(v) \subseteq V_{\ell+1}$ drawn from the already embedded coarser level.
The initial location of $v$ is obtained by minimizing
\[
  \Phi_\ell(\bz_v)
  =
  \sum_{a \in A_\ell(v)}
  \omega_{va}
  \bigl( \norm{\bz_v - \bz_a} - d_G(v,a) \bigr)^2,
\]
where the anchor weights $\omega_{va}$ favor well-spread, geodesically close
anchors.  This is a local nonlinear least-squares problem and plays the role of
geodesic trilateration.  In the current implementation each inserted vertex is
optimized independently by a small Armijo line-search solver.  The middle
panels in the lower rows of
\cref{fig:misf_stages_regular_12x12,fig:misf_stages_regular_15x15,fig:misf_stages_irregular_15x15}
show these insertion states before any active-set refinement.

\paragraph{Sparse active-set refinement.}
After all vertices of level $V_\ell$ have been inserted, the active set is
refined by a sparse GMDS solve.  Let $P_\ell \subseteq V_\ell \times V_\ell$ be
a sparse pair set consisting of short-range geodesic neighbors together with a
small number of landmark pairs.  The level objective has the form
\[
  F_\ell(\bZ)
  =
  \sum_{(i,j)\in P_\ell}
  \bigl(
    \demb_{\Gamma}(i,j;\bZ) - d_G(i,j)
  \bigr)^2
  +
  \lambda_\ell
  \sum_{u \notin V_\ell}
  \alpha_u \norm{\bz_u - \bar{\bz}_u}^2,
\]
where $\bar{\bz}_u$ denotes the current scaffold position of an inactive
vertex.  The anchor term pins the already constructed coarser scaffold while
allowing the active vertices to adjust.  In other words, the method uses the
same fixed-geodesic discrepancy as GMDS, but only on a sparse set of active
pairs, and with the non-active scaffold softly constrained.  The right panels
of \cref{fig:misf_stages_regular_12x12,fig:misf_stages_regular_15x15,fig:misf_stages_irregular_15x15}
show the corresponding refinement states.  After the levelwise cycle reaches
$V_0$, an optional short full-graph pure-GMDS polish is applied.

\paragraph{Implementation in \texttt{grip}.}
The current prototype reuses the GRIP infrastructure already present in the
\texttt{grip} package.  The MISF itself is extracted by
\texttt{grip.build.misf()} in \texttt{R/grip\_misf.R}, which delegates to the
compiled GRIP filtration code in \texttt{src/DrawGraph.cpp} and
\texttt{src/grip\_rcpp.cpp}.  The multiscale GMDS preparation and orchestration
are implemented in \texttt{R/grip\_geodesic\_misf.R} through
\texttt{grip.prepare.misf.geodesic.mds()} and
\texttt{grip.optimize.misf.geodesic.mds()}.  The per-vertex insertion kernel is
implemented in \texttt{src/geodesic\_misf\_rcpp.cpp}, while both the coarse
top-level solve and the sparse refinement stages reuse the compiled GMDS kernel
in \texttt{src/geodesic\_mds\_rcpp.cpp}.  Thus the MISF modification changes the
optimization schedule and the initialization strategy, but not the underlying
fixed-path geodesic discrepancy.

\paragraph{What the multiscale figures show.}
\Cref{fig:misf_stages_regular_12x12} shows the regular $12\times 12$
paraboloid.  The coarsest initial embedding on $V_2$ is already far from a
recognizable paraboloid sample ($\rho = 0.8588$), and the subsequent
refinement improves it only to $\rho = 0.7109$.  At the full level $V_0$, the
insertion state is highly tangled and refinement reduces the active-set error
from $\rho = 0.6772$ to $\rho = 0.5763$, but the result is still not a smooth
paraboloid.  \Cref{fig:misf_stages_regular_15x15} shows the same phenomenon on
the regular $15\times 15$ mesh: every level improves under refinement, yet the
coarse sample and the full-level embedding remain visibly far from the target
surface.

By contrast, \cref{fig:misf_stages_irregular_15x15} shows that the irregular
rectangle paraboloid is substantially easier for the same multiscale pipeline.
Its coarsest initial embedding starts much closer to the target
($\rho = 0.5331$), and each subsequent level remains more consistent with the
reference surface.  The final level improves from $\rho = 0.3617$ after
insertion to $\rho = 0.3343$ after sparse refinement, followed by a final full
polish at $\rho = 0.3205$.  This mirrors the empirical observation from the
previous section: the difficulty is not merely the presence of curvature, but
the interaction between the optimization landscape and the high symmetry of the
regular rectangular mesh.

\begin{table}[t]
  \centering
  \small
  \caption{Stagewise active-set Procrustes error $\rho$ for the MISF-GMDS
  pipeline on the three paraboloid families shown in
  \cref{fig:misf_stages_regular_12x12,fig:misf_stages_regular_15x15,fig:misf_stages_irregular_15x15}.  The
  ``initial'' column is the coarse top-level GMDS solve for the first active
  level and the anchor-based insertion state for the remaining levels; the
  ``refined'' column is the corresponding sparse GMDS refinement.}
  \label{tab:misf_stage_rho}
  \begin{tabular}{llrrr}
    \toprule
    case & level & $|V_\ell|$ & initial $\rho$ & refined $\rho$ \\
    \midrule
    regular $12\times 12$ & $V_2$ & 24  & 0.8588 & 0.7109 \\
    regular $12\times 12$ & $V_1$ & 50  & 0.8226 & 0.6731 \\
    regular $12\times 12$ & $V_0$ & 144 & 0.6772 & 0.5763 \\
    regular $12\times 12$ & final & 144 & ---    & 0.5792 \\
    \midrule
    regular $15\times 15$ & $V_3$ & 24  & 0.9181 & 0.7114 \\
    regular $15\times 15$ & $V_2$ & 28  & 0.7335 & 0.6137 \\
    regular $15\times 15$ & $V_1$ & 84  & 0.6077 & 0.4897 \\
    regular $15\times 15$ & $V_0$ & 225 & 0.5709 & 0.4751 \\
    regular $15\times 15$ & final & 225 & ---    & 0.4754 \\
    \midrule
    irregular $15\times 15$ & $V_3$ & 24  & 0.5331 & 0.4684 \\
    irregular $15\times 15$ & $V_2$ & 32  & 0.4970 & 0.4221 \\
    irregular $15\times 15$ & $V_1$ & 81  & 0.3969 & 0.3555 \\
    irregular $15\times 15$ & $V_0$ & 225 & 0.3617 & 0.3343 \\
    irregular $15\times 15$ & final & 225 & ---    & 0.3205 \\
    \bottomrule
  \end{tabular}
\end{table}

\begin{figure}[t]
  \centering
  \includegraphics[width=\linewidth]{../figures/fig_misf_filtration_cases.pdf}
  \caption{MISF levels displayed on the underlying $2$D mesh geometry.  Each
  row shows one graph family, with the active level $V_\ell$ highlighted on top
  of the full mesh graph.  The irregular rectangle keeps rectangular topology
  but breaks the lattice symmetry in the geometry, yielding a more generic
  coarse sample than the perfectly regular meshes.}
  \label{fig:misf_filtration_cases}
\end{figure}

\begin{figure}[t]
  \centering
  \includegraphics[width=\linewidth]{../figures/fig_misf_stages_regular_12x12.pdf}
  \caption{Regular $12\times 12$ paraboloid under MISF-GMDS.  Each row
  corresponds to one active set $V_\ell$.  Left: the reference active-set
  sample.  Middle: the initial multiscale embedding for that level (the coarse
  pure-GMDS solve for $V_2$, then anchor-based insertion for $V_1$ and $V_0$).
  Right: the sparse GMDS refinement of the same active set.  The coarsest level
  is already not paraboloid-like, and although every level improves, the full
  embedding remains folded.}
  \label{fig:misf_stages_regular_12x12}
\end{figure}

\begin{figure}[t]
  \centering
  \includegraphics[width=\linewidth]{../figures/fig_misf_stages_regular_15x15.pdf}
  \caption{Regular $15\times 15$ paraboloid under MISF-GMDS.  The same
  qualitative pattern persists at the larger size: the coarsest-level solve does
  not recover a recognizable paraboloid sample, insertion creates a tangled
  full-level state, and sparse refinement improves but does not remove the
  pathology.}
  \label{fig:misf_stages_regular_15x15}
\end{figure}

\begin{figure}[t]
  \centering
  \includegraphics[width=\linewidth]{../figures/fig_misf_stages_irregular_rectangle_15x15.pdf}
  \caption{Irregular $15\times 15$ rectangle lifted to a paraboloid.  The
  multiscale pipeline is much better behaved on this less symmetric graph: the
  coarsest level already resembles a sparse paraboloid sample, and each finer
  refinement remains visibly closer to the reference geometry than in the
  regular-grid cases.}
  \label{fig:misf_stages_irregular_15x15}
\end{figure}

\section{Experiments: What the Current Implementation Teaches Us}
\label{sec:experiments_v4}

The goal of the current experimental section is not to present a finished
method, but to identify which part of the working implementation is
responsible for the main pathologies and which interventions already produce
meaningful improvements.

\subsection{Current State of the Working Implementation}
\label{sec:current_implementation}

The working implementation in the \texttt{grip} package is now best viewed as
a family of related GMDS prototypes rather than a single finished algorithm.
The common core is still the fixed-geodesic discrepancy
\cref{eq:raw_geo_stress}: for a prepared graph-geodesic family
$(G,\Gamma)$, the implementation measures the embedded length of the
input-side geodesics and optimizes that mismatch in a compiled C++ kernel.
Around that core, several experimental layers have been added in order to
understand the main failure modes of the objective.

\begin{table}[t]
  \centering
  \small
  \caption{Current GMDS implementation layers in \texttt{grip}.}
  \label{tab:implementation_layers}
  \begin{tabular}{p{0.23\linewidth} p{0.34\linewidth} p{0.35\linewidth}}
    \toprule
    layer & main files / entry points & current role \\
    \midrule
    Pure and regularized GMDS &
    \texttt{R/grip\_quality.R},
    \texttt{src/geodesic\_mds\_rcpp.cpp},
    \texttt{grip.prepare.graph.geodesic.mds()},
    \texttt{grip.optimize.geodesic.mds()} &
    Fixed-geodesic objective, tie-averaged shortest paths, anchor regularization, graph-distance-aware repulsion, and edge-spring relaxations. \\
    \addlinespace
    MISF-based GMDS &
    \texttt{R/grip\_misf.R},
    \texttt{R/grip\_geodesic\_misf.R},
    \texttt{src/geodesic\_misf\_rcpp.cpp},
    \texttt{grip.optimize.misf.geodesic.mds()} &
    Coarse-to-fine GMDS pipeline with top-level solve, per-vertex insertion, sparse level refinement, and final full-graph polish. \\
    \addlinespace
    Layout-based basin finders &
    \texttt{R/grip\_layout\_weighted.R},
    \texttt{src/MishWeighted.cpp},
    \texttt{src/grip\_rcpp.cpp},
    Weighted GRIP / weighted KK / LGKK APIs &
    External basin-finding and warm-start machinery for MISF-GMDS; currently the most useful source of coarse scaffolds and lower-level placement states. \\
    \bottomrule
  \end{tabular}
\end{table}

At this stage the implementation already supports several experimentally
distinct variants: pure GMDS; tie-averaged GMDS; anchor-regularized GMDS;
anchor-plus-repulsion and anchor-plus-edge-spring regularizations; an
edge-relaxation surrogate; a MISF-based multiscale GMDS schedule; and seeded
MISF variants that reuse KK, GRIP, and Weighted GRIP layouts as basin
finders.  The goal of the recent experimental cycle was not merely to optimize
one more hyperparameter, but to determine \emph{which part of the current
pipeline is actually responsible for the observed pathologies}.

\subsection{What the Recent Experiments Show}
\label{sec:recent_experiments}

\paragraph{The smooth target may exist even when the optimizer misses it.}
The identifiability experiment on the regular and irregular paraboloid meshes
showed that the smooth paraboloid is an \emph{exact} GMDS solution after the
graph-weight normalization used by the generator.  In particular, for the
regular $12\times 12$ paraboloid, a $15^\circ$ folded perturbation of the
truth has only moderate raw geodesic stress
($\sigma = 0.0102$, $\rho = 0.0325$), and pure GMDS started from that folded
state returns close to the smooth reference
($\sigma = 0.0002$, $\rho = 0.0303$ after $25$ iterations).  By contrast,
pure GMDS started from a global classical-MDS seed lands in a very different
basin ($\sigma = 0.0298$, $\rho = 0.3291$ on the same graph).  Thus the main
issue is not absence of a smooth solution but basin selection.

\paragraph{Single-scale regularization helps locally but does not solve the
basin problem.}
On the regular paraboloid meshes, anchor-based regularization improves the raw
geodesic objective substantially but leaves the recovered shape much too far
from the target.  For the regular $12\times 12$ paraboloid, direct
\texttt{cmdscale $\rightarrow$ pure GMDS} ends at
$(\sigma,\rho)=(0.0492, 0.3257)$, while the best single-scale regularized run
reaches $(0.0221, 0.3214)$.  For the regular $15\times 15$ paraboloid, the
best regularized run reaches $(0.0183, 0.2286)$, improving stress much more
than shape.  In other words, the current anchor, repulsion, and edge-spring
terms are useful as controlled perturbations of the objective, but they do not
yet provide a reliable route to the correct basin on their own.

\paragraph{Multiscale GMDS helps only if the scaffold is right.}
The first MISF-GMDS prototype in \cref{sec:misf_gmds} showed that a pure-GMDS
coarse solve is often not a good top-level scaffold on regular meshes.  The
subsequent Phase~A--E experiments therefore replaced ``one GMDS initializer''
by a seed pool and asked which part of the multiscale schedule actually needs
help.  The answer is now fairly clear:
\begin{enumerate}
  \item the top MISF level is a basin-selection problem, not merely a local
  refinement problem;
  \item lower-level geodesic insertion can destroy a good coarse scaffold;
  \item weighted KK is currently the strongest lower-level placement mechanism;
  \item a small pool of top-level seeds is better than any single universal
  default.
\end{enumerate}

\begin{table}[t]
  \centering
  \small
  \caption{Representative recent outcomes on the paraboloid families.  The
  ``direct'' column is the full-graph \texttt{cmdscale $\rightarrow$ pure GMDS}
  control.  The ``best single-scale regularized'' column reports the best
  regularized-GMDS run from the single-scale study.  The last column gives the
  best currently observed seeded multiscale pipeline on the corresponding
  family.}
  \label{tab:recent_paraboloid_results}
  \begin{tabular}{lccc}
    \toprule
    case & direct $(\sigma,\rho)$ & best single-scale regularized $(\sigma,\rho)$ & current best multiscale $(\sigma,\rho)$ \\
    \midrule
    regular paraboloid $12\times 12$ & $(0.0492, 0.3257)$ & $(0.0221, 0.3214)$ & $(0.0109, 0.0904)$ \\
    regular paraboloid $15\times 15$ & $(0.0519, 0.2317)$ & $(0.0183, 0.2286)$ & $(0.0148, 0.0726)$ \\
    irregular rectangle paraboloid $15\times 15$ & $(0.2872, 0.1845)$ & --- & $(0.0951, 0.0889)$ \\
    \bottomrule
  \end{tabular}
\end{table}

\begin{figure}[t]
  \centering
  \includegraphics[width=0.98\linewidth]{../figures/fig_current_implementation_summary.pdf}
  \caption{Summary of the current implementation state.  Left: progression of
  final Procrustes error $\rho$ on the three paraboloid families as the
  pipeline moves from direct full-graph GMDS to seeded multiscale variants.
  Right: Block~1 of the Phase~E mesh-geometry study shows that the current
  truth-free selector is still too \texttt{cmdscale}-biased: it chose
  \texttt{cmdscale} in $10/12$ cases, whereas the oracle best final-$\rho$
  seed was \texttt{kk} in $6/12$ cases, \texttt{cmdscale} in only $3/12$,
  \texttt{grip} in $2/12$, and \texttt{weighted\_grip\_polish\_lgkk} in
  $1/12$.}
  \label{fig:current_implementation_summary}
\end{figure}

\paragraph{The current selector is not yet good enough.}
The most recent Block~1 mesh-geometry expansion confirms that there is no
single universal top-level winner across flat, saddle, paraboloid, and ripple
meshes.  The current truth-free proxy chose \texttt{cmdscale} in $10/12$
cases and matched the oracle best final-$\rho$ seed in only $4/12$ cases.
This is precisely why the current pipeline should still be viewed as a
research prototype: the best-performing top seed depends on geometry, and the
automatic selector still over-rewards immediate low-stress spectral states.

\section{Discussion and Immediate Next Steps}
\label{sec:discussion_v4}

The current state of the project is therefore scientifically useful even
though the implementation is still incomplete: the theory is clearer, the
optimizer failure modes are better localized, and the remaining work can now
be organized as a targeted experimental program rather than a broad search.

\subsection{Immediate Next Steps}
\label{sec:next_steps}

The recent experiments reduce the open problem to a much narrower one than we
had at the beginning of the project.  The main challenge is no longer to
decide whether GMDS should be single-scale or multiscale, or whether classical
MDS should be discarded entirely.  Instead, the present evidence suggests the
following agenda.

\paragraph{1. Improve the truth-free seed selector.}
The current Phase~E proxy uses only a one-level weighted-KK lookahead together
with three rank terms (lookahead stress, roughness, and edge-floor).  This is
not yet predictive enough.  The next selector experiments should include
insertion-stability signals such as area-floor diagnostics, first-level energy
jumps, and possibly a second short refinement step before choosing the winner.
An especially promising direction is to retain a small shortlist of seeds
rather than forcing a single winner too early.

\paragraph{2. Keep weighted KK as the current lower-level default.}
Across the recent multiscale studies, weighted KK is the strongest lower-level
placement rule we currently have.  It should therefore remain the default
lower-level mechanism while the top-level seed-selection problem is being
improved.  Future comparisons should treat GRIP-family lower-level placement as
an ablation against this weighted-KK baseline, not as the default.

\paragraph{3. Expand only after the selector improves.}
The ongoing mesh-geometry expansion is useful because it reveals where the
selector fails, but the next full cross-family expansion should wait until the
top-level proxy is better calibrated.  Otherwise we risk spending substantial
computation on a pipeline whose automatic choice rule is not yet trustworthy.
The most natural order is: first improve the selector on the mesh geometries,
then move to perforated rectangles, and only then to annulus/sphere families.

\paragraph{4. Continue objective-level diagnostics without conflating them with
the optimizer.}
The experiments already show that the smooth target may be an exact GMDS
solution while the current initialization misses it.  This means that
objective-level questions (fold-blindness, path averaging, repulsion,
edge-relaxation) and optimization-level questions (basin finding, multiscale
placement, seed selection) should continue to be studied in parallel but kept
conceptually separate.  At present, the most urgent unresolved issue is
optimization/basin selection rather than the formal definition of GMDS itself.

\appendix

\section{Background on MDS, Isomap, and Related Methods}
\label{app:mds_background}

This appendix records background material that supports the main text but is
not essential to the primary narrative.

\subsection{SMACOF and Majorization}
\label{sec:smacof}

Direct minimization of weighted Euclidean stress is non-convex and typically
handled by SMACOF \citep{deleeuw1977applications,deleeuw2009multidimensional}.
The key idea is to majorize the stress by a quadratic surrogate that is tight
at the current iterate, leading to a monotone descent scheme.  This classical
viewpoint is useful because it motivates the path-majorization heuristics used
later for GMDS, even though the geodesic objective is more complicated than the
standard Euclidean one.

\subsection{Non-Metric MDS}
\label{sec:nonmetric_mds}

Non-metric MDS \citep{kruskal1964nonmetric,shepard1962analysis} relaxes the
requirement to match numerical dissimilarities exactly and instead preserves
their rank order through isotonic regression.  A non-metric geodesic variant of
GMDS would apply the same idea to graph-geodesic targets while continuing to
measure output fidelity by embedded path lengths rather than by Euclidean
distances.

\subsection{Related Manifold-Learning and Graph-Drawing Methods}
\label{sec:manifold_learning}

The broader context includes manifold-learning methods such as LLE
\citep{roweis2000nonlinear}, Laplacian Eigenmaps \citep{belkin2003laplacian},
Hessian Eigenmaps \citep{donoho2003hessian}, Diffusion Maps
\citep{coifman2006diffusion}, $t$-SNE \citep{vandermaaten2008tsne}, and UMAP
\citep{mcinnes2018umap}.  Most of these methods still compare input structure
with Euclidean output distances.  Graph-drawing methods such as Kamada-Kawai
\citep{kamada1989algorithm} and Fruchterman-Reingold
\citep{fruchterman1991graph} are also relevant because they provide practical
stress-based or force-based layout seeds that turn out to be important for the
working GMDS implementations discussed in the paper.

\section{Gradient and Majorization Details}
\label{app:gradient_majorization}

This appendix records the detailed gradient structure of fixed-path GMDS and a
majorization-style optimization viewpoint.

\subsection{Path-Nonlocal Gradient}

For each pair $i<j$, let
$\gamma_{ij}^{(k)}=(v_0=i,v_1,\ldots,v_L=j)$ be the chosen shortest path.
Define the unit edge vector
\begin{equation}
  \mathbf{e}_{ab} = \frac{\bz_a - \bz_b}{\norm{\bz_a - \bz_b}},
  \label{eq:unit_edge}
\end{equation}
away from degenerate zero-length edges.  The derivative of the embedded path
length with respect to a vertex $\bz_\ell$ is
\begin{align}
  \frac{\partial\, \demb_{\Gamma_k}(i,j;\bZ)}{\partial\, \bz_\ell}
  &= \begin{cases}
       \mathbf{e}_{\ell, v_{s-1}} + \mathbf{e}_{\ell, v_{s+1}}
         & \text{if } \ell = v_s,\ 0 < s < L, \\
       \mathbf{e}_{\ell, v_1}
         & \text{if } \ell = v_0 = i, \\
       \mathbf{e}_{\ell, v_{L-1}}
         & \text{if } \ell = v_L = j, \\
       \bm{0}
         & \text{if } \ell \notin \gamma_{ij}^{(k)}.
     \end{cases}
  \label{eq:geo_distance_gradient}
\end{align}
Consequently,
\begin{equation}
  \frac{\partial\, \sigma^{\mathrm{geo}}}{\partial\, \bz_\ell}
  = 2 \sum_{i < j}
    \bigl(\demb_{\Gamma_k}(i,j;\bZ) - \delta^{(k)}_{ij}\bigr)
    \cdot \frac{\partial\, \demb_{\Gamma_k}(i,j;\bZ)}{\partial\, \bz_\ell}.
  \label{eq:raw_gradient}
\end{equation}
The important structural fact is that the gradient at $\ell$ receives
contributions from every endpoint pair whose chosen input geodesic passes
through $\ell$.  GMDS is therefore non-local even when the underlying graph is
sparse.

\subsection{A Majorization Viewpoint}

The standard SMACOF majorization exploits the fact that Euclidean distance is a
single norm.  In GMDS, the output-side discrepancy is a sum of norms along a
fixed path.  If $\bZ^{(t)}$ is the current iterate, then each edge norm admits
the Cauchy-Schwarz majorization
\begin{equation}
  \norm{\bz_a - \bz_b}
  \leq
  \frac{\norm{\bz_a - \bz_b}^2 + \norm{\bz^{(t)}_a - \bz^{(t)}_b}^2}
       {2\,\norm{\bz^{(t)}_a - \bz^{(t)}_b}},
  \label{eq:cs_majorization}
\end{equation}
which can be accumulated over the stored geodesic family.  This motivates the
path-edge weights
\begin{equation}
  \tilde{w}_{ab}^{(t)}
  = \sum_{i < j : (a,b) \in \gamma_{ij}^{(k)}}
    \frac{1}{\norm{\bz^{(t)}_a - \bz^{(t)}_b}},
  \label{eq:path_edge_weight}
\end{equation}
and the corresponding path-augmented Guttman matrix
\begin{equation}
  \tilde{\bV}^{(t)}_{ab}
  = -\tilde{w}^{(t)}_{ab}, \quad a \neq b; \qquad
  \tilde{\bV}^{(t)}_{aa}
  = \sum_{b \neq a} \tilde{w}^{(t)}_{ab}.
  \label{eq:augmented_guttman}
\end{equation}
A natural update is then
\begin{equation}
  \bZ^{(t+1)} = \bigl(\tilde{\bW}^{(t)}\bigr)^+ \tilde{\bV}^{(t)} \bZ^{(t)},
  \label{eq:geo_guttman_update}
\end{equation}
where $\tilde{\bW}^{(t)}$ is the diagonal matrix of row sums of
$|\tilde{\bV}^{(t)}|$.

\begin{proposition}[Monotone decrease under fixed-path majorization]
  \label{prop:monotone}
  With the geodesic family fixed, a valid majorization step of the form
  \eqref{eq:geo_guttman_update} does not increase the corresponding surrogate
  objective.
\end{proposition}

\begin{proof}[Proof sketch]
  The pathwise majorization produces an upper bound that is tight at the current
  iterate.  Minimizing that upper bound yields a monotone descent step for the
  surrogate.  The practical importance of this viewpoint is less the exact
  formula than the observation that GMDS admits structured iterative solvers
  that reuse the same cached path family throughout optimization.
\end{proof}

\section{Metric-Geometric Context}
\label{sec:metric_geometry_v4}
\label{sec:metric_geometry}

This appendix records the metric-geometric ideas that help interpret GMDS but
are not part of the primary narrative.

\paragraph{Minimum realization.}
The $k$-NN graph $\Gk$, after removing metrically redundant edges, yields the
\emph{minimum realization} of the input graph metric
\citep{dress1984trees, imrich1984optimal}.  GMDS on that realization uses the
most parsimonious graph structure compatible with the chosen metric.

\paragraph{Tight span.}
The \emph{tight span} $T(\Delta)$ (injective hull) \citep{dress1984trees,
isbell1964six} is the minimal injective metric space containing the finite
metric $(X,\Delta)$.  GMDS may be viewed as trying to preserve, in Euclidean
coordinates, the path structure already encoded in the one-skeleton of that
metric realization.

\paragraph{Metric spanners.}
Neighborhood graphs can often be interpreted as sparse spanners of the ambient
metric \citep{althofer1993sparse, har-peled2006low}.  From that perspective,
GMDS trades dense Euclidean pair information for a sparse shortest-path model
that preserves geometry only up to a chosen scale.

\paragraph{Gromov-Hausdorff approximation.}
Under suitable sampling conditions, graph shortest-path metrics converge toward
the underlying geodesic metric space in the Gromov-Hausdorff sense
\citep{gromov1999metric, latschev2001riemannian, bernstein2000graph}.  This
gives one metric-geometric interpretation of the data-modeling step discussed in
the main text.

%% ============================================================
%% BIBLIOGRAPHY
%% ============================================================

\section*{Acknowledgments}
Research reported in this publication was supported by grant number
INV-048956 from the Gates Foundation.

\bibliographystyle{plainnat}
\bibliography{geodesic_mds}

\end{document}
